
\documentclass{article}
\usepackage[final]{neurips}
\usepackage{graphicx} % Required for inserting images
\usepackage{natbib}
\usepackage[utf8]{inputenc} % allow utf-8 input
\usepackage[T1]{fontenc}    % use 8-bit T1 fonts
\usepackage{hyperref}       % hyperlinks
\usepackage{url}            % simple URL typesetting
\usepackage{booktabs}       % professional-quality tables
\usepackage{amsfonts}       % blackboard math symbols
\usepackage{nicefrac}       % compact symbols for 1/2, etc.
\usepackage{microtype}      % microtypography
\usepackage{xcolor}         % colors
\usepackage{ifdraft}

\usepackage{amsmath}
\usepackage{graphicx}
\usepackage{float}
\usepackage{subcaption}
\usepackage{enumitem}
\usepackage{fancyvrb}
\usepackage{amssymb}
\usepackage{cleveref}
\usepackage{tcolorbox}
\usepackage{newtxtext}
% \usepackage{libertine} %this is breaking my verbatim.
\usepackage{soul}

\usepackage{pifont}

\newcommand{\yes}{\ding{51}} % ✓
\newcommand{\no}{\ding{55}} 

\title{Auto-formalizing Research Mathematics into Verified Lean Proofs}
\author{Raymond Xu, Yuchen Hou\\
  {\small \texttt{xu.raym@northeastern.edu, hou.yuchen@northeastern.edu}}}
\date{July 2026}

\begin{document}

\maketitle

\section{Introduction}
LLM scalability has three dimensions:  training, inference and evaluation. Training scalability concerns efficiently using computation to train larger models; inference scalability concerns allocating test-time computation to improve solutions; and evaluation scalability concerns verifying those solutions with minimal human effort. In software, compilers provide an automatic correctness signal. In mathematics, proof assistants such as Lean can play a similar role by checking whether a solution is formally valid.

However, most mathematical problems and proofs are written in natural language, while formal verification requires expert-written formal statements and proofs. This manual translation creates a major bottleneck. Autoformalization—automatically translating natural-language mathematics into verifiable formal code—is therefore essential for scaling the reliable evaluation of LLM reasoning. 

We study \emph{proof autoformalization}: translating a natural-language theorem
together with its proof, $(NL_{\text{thm}}, NL_{\text{proof}})$, into a Lean~4
statement and proof, $(FL_{\text{thm}}, FL_{\text{proof}})$, that the kernel
accepts and that follows the same argument~\citep{lean4}. Recent systems show
the task is tractable at competition and undergraduate level: ProofBridge
translates full theorem-and-proof pairs with joint NL--FL embeddings and
iterative repair~\citep{proofbridge2025}, and ProofFlow decomposes an informal
proof into a dependency graph of lemma statements and formalizes each
node~\citep{proofflow2025}. At research level the numbers stay low: the best
reported pass rate on RLMEval, built from real Lean blueprint projects, is
10.3\%~\citep{rlmeval2025}, and open provers solve 3\% of FATE-H and none of
FATE-X~\citep{fate2025}. We target this gap with three questions: how to ground
generated code in the actual Mathlib namespace, where structure-guided
translation fails once grounding is controlled, and how to train small open
models with a dense and semantically meaningful verifier signal.

Our contributions are: (1) a grounding mechanism that constrains identifier
decoding to the declaration names of a pinned Mathlib build, removing the
unknown-identifier failure class at the statement level by construction; (2) a
preregistered factorial study on research-level benchmarks that separates the
contributions of structure, retrieval, and grounding, and labels every failure
into five classes; (3) a per-node reward for reinforcement learning on proof
translation that combines node compilation, node provability, and node-level
semantic consistency with the aligned informal step, compared against
outcome-only and syntactic per-node rewards under matched budgets; and (4) a
soundness-audited evaluation harness with Comparator statement matching, axiom
checks, and per-problem cost reporting. Two observations motivate the design.

%We want to build a system that turns a math problem into a verifiable Lean~\citep{lean4}
%proof that the Lean kernel actually accepts. This happens in two stages:
%\emph{autoformalization} (natural language $\rightarrow$ Lean statement) and
%\emph{proof generation} (Lean statement $\rightarrow$ Lean proof). Two
%observations motivate the project.

\paragraph{Library grounding gates small-model formalization.} At the frontier,
surface syntax has stopped being the binding constraint: agentic pipelines emit
almost uniformly parseable Lean, and competition benchmarks are close to
saturation, with Seed-Prover at 99.6\% on MiniF2F~\citep{seedprover2025} and
Goedel-Architect closing the full test split with natural-language
guidance~\citep{goedelarchitect2026}. The constraint that persists across scales
is grounding in the library and fidelity of the stated claim. A frontier model
still produces unknown identifiers on roughly half of ProofNet\# items under
paraphrase~\citep{surfacesens2026}, agentic formalizers reach 89.5\% compilation
while 60.5\% of outputs are faithful to the source
statement~\citep{beyondcompilation2026}, and research-level benchmarks remain in
the single digits~\citep{rlmeval2025, fate2025}. Our own failure analysis in
Section~\ref{sec:results} shows the same shift for the 7B class we study:
finetuning repairs surface syntax and moves the dominant failure to hallucinated
library names. We therefore treat grounding and faithfulness as the quantities
to improve.

%\paragraph{Formalization ability gates proof ability.} How well a model proves
%depends on how well it formalizes. This is obvious at stage one, but it also
%shows up at stage two: proof generators often produce proofs that simply don't
%compile because they get Lean syntax and library usage wrong. So we treat
%formalization skill as the main thing to improve, not just a preprocessing step.

\paragraph{Aligned proof pairs exist and are underused.} Reference proofs are
plentiful in synthetic form, though scarce at research level. Herald
informalizes Mathlib4 and releases 44k aligned NL--FL proof pairs next to 580k
statement pairs~\citep{herald2025}, and ProofFlow provides 184 undergraduate
problems with annotated step dependency graphs~\citep{proofflow2025}. Existing
pipelines use this material mainly to train statement formalizers. We use the
proof pairs directly as supervision for translation, restructured into the
skeleton format of Section~\ref{sec:method}, and we reserve blueprint-project
alignments~\citep{rlmeval2025} for research-level evaluation, giving the model a
denser training signal than the compiler's 0/1 reward.

%\paragraph{Ground-truth proofs are scarce.} Most Lean benchmarks give statements
%but no reference proofs. That leaves the proof generator to learn from a sparse
%0/1 signal: does it compile, and does the Lean comparator confirm it proves the
%right theorem. At the same time, lots of proof material exists in natural
%language and goes unused. We want to use both natural-language proofs as
%extra supervision, and any available ground-truth values to give the model a
%denser training signal than the compiler's 0/1 reward.

\section{Related Work}
{\color{red}
\paragraph{Statement autoformalization.} LeanEuclid targets Euclidean
geometry~\citep{leaneuclid2024}. Sampling with type-check filtering and
Self-BLEU selection reaches 45.1\% on ProofNet\# and 31.6 BEq+ on RLM25 with
GPT-4o at large budgets~\citep{poiroux2025beq}, and RAutoformalizer adds
dependency retrieval with a faithfulness-oriented metric, improving BEq@8 from
12.83\% to 18.18\% on ProofNet~\citep{rautoformalizer2025}. Reinforcement
learning now leads this subtask: FormaRL lifts Qwen2.5-Coder-7B from 4.04\% to
26.15\% ProofNet pass@1 with 859 unlabeled examples using a compiler reward and
an LLM consistency reward~\citep{formarl2025}; Mathesis trains the formalizer
online with GRPO and prover-preference optimization, introducing the LeanScorer
semantic metric and the Gaokao-Formal benchmark~\citep{mathesis2025}; CriticLean
trains the critic itself and releases the 285k-statement
FineLeanCorpus~\citep{criticlean2025}. Herald builds 580k statement pairs and
44k proof pairs by informalizing Mathlib4~\citep{herald2025}. Two recent
analyses shape our design: paraphrase probes find unknown identifiers on about
half of ProofNet\# items for a strong closed model and name constrained decoding
over the Mathlib interface as an open direction~\citep{surfacesens2026}, and a
factorial evaluation of agentic formalization reports a 29-point gap between
compilation rate and faithfulness~\citep{beyondcompilation2026}.

\paragraph{Proof generation.} ReProver couples retrieval over accessible
premises with tactic generation~\citep{leandojo2023}. Provers trained with
reinforcement learning have since saturated competition benchmarks:
DeepSeek-Prover-V2 trains on subgoal decomposition~\citep{deepseekproverv2},
Goedel-Prover-V2 adds scaffolded data synthesis and compiler-guided
self-correction, reaching 88.1\% pass@32 on MiniF2F with a 32B
model~\citep{goedelproverv2_2025}, Kimina-Prover scales reasoning-style proving
to 72B~\citep{kimina_prover_2025}, and Seed-Prover reaches 99.6\% on MiniF2F and
5 of 6 problems at IMO 2025~\citep{seedprover2025}, with its successor at 87.9\%
on PutnamBench~\citep{seedprover15_2025}. Goedel-Architect organizes proving
around a blueprint, a dependency graph of definitions and lemmas whose nodes are
closed in parallel and refined on failure, at a reported cost of \$0.44 per
problem~\citep{goedelarchitect2026, leanarchitect2026}. The line descends from
Draft, Sketch, and Prove, which seeded formal sketches with informal
proofs~\citep{dsp2023}.

\paragraph{Proof-level autoformalization.} Translating the proof itself has two
recent treatments. ProofBridge trains a joint embedding over NL--FL
theorem-and-proof pairs, retrieves aligned exemplars, and repairs iteratively
with type-checker and equivalence feedback, improving semantic correctness by
31.14 points over Kimina-Prover-RL-1.7B at pass@32 on a curated MiniF2F
proof-pair split~\citep{proofbridge2025}. ProofFlow maps an informal proof to a
dependency graph of intermediate lemmas with \texttt{sorry} placeholders and
reports a ProofScore of 0.545 against 0.123 for monolithic full-proof
formalization~\citep{proofflow2025}. Both operate at competition and
undergraduate difficulty. Our work differs in three ways: we ground identifier
generation in the library by construction, we measure where structure-guided
translation fails at research level with a labeled failure taxonomy, and we
train the translator with a per-node semantic reward.

\paragraph{Training signal and evaluation.} Dense verifier feedback for proving
is established: tactic-level rewards from Lean elaboration with first-error
credit outperform outcome-only GRPO on MiniF2F and
ProofNet~\citep{processrl2026}, and FoVer trains process reward models on step
labels from Z3 and Isabelle~\citep{fover2025}. Semantic rewards exist so far
only at the statement level~\citep{formarl2025, mathesis2025, criticlean2025}.
On evaluation, an audit of Lean benchmarks documents specification errors and
scoring loopholes, showing that kernel acceptance certifies a narrower claim
than users often assume~\citep{faultsbench2026}; FormalQualBench operationalizes
specification-based scoring with the Lean FRO Comparator on 23 qualifying-exam
theorems, where the strongest published harnesses solve 8 and 10
problems~\citep{formalqualbench2026, opengauss2026, merlean2026}. We adopt this
standard together with the unbiased pass@$k$ estimator~\citep{chen2021codex}.
}

% \paragraph{Baselines.} On \emph{statement autoformalization} (NL statement
% $\rightarrow$ Lean statement), LeanEuclid~\citep{leaneuclid2024} is a baseline in
% the domain of Euclidean geometry. A separate line of work uses 12-shot prompting
% with type-check filtering and Self-BLEU selection, generating many candidates
% and picking one with a heuristic, reaching 45.1\% on ProofNet\# and 31.6 BEq+
% on RLM25 with GPT-4o, but at large sampling budgets.
% RAutoformalizer~\citep{rautoformalizer2025} retrieves library dependencies and
% adds a faithfulness-oriented metric, improving BEq@8 from 12.83\% to 18.18\% on
% ProofNet, though it requires dependency retrieval over Lean library context.
% FormaRL~\citep{formarl2025} uses RL from
% Lean syntax checks plus an LLM semantic-consistency reward, and is data-efficient:
% with 859 unlabeled examples it lifts Qwen2.5-Coder-7B ProofNet pass@1 from 4.04\%
% to 26.15\%. On \emph{proof generation} (formal statement $\rightarrow$ Lean
% proof), LeanDojo's ReProver~\citep{leandojo2023} is a retrieval-augmented prover
% over accessible premises.

\paragraph{Benchmarks.} Table~\ref{tab:benchmarks} lists the main ones. Two
things stand out: most pair NL and formal statements but include no reference
proofs (the sparse-reward setting above), and difficulty ranges from olympiad
(MiniF2F, PutnamBench) up to graduate and research level (RLMEval, FATE,
FormalQualBench), where the strongest systems remain far from saturation.

\begin{table}[H]
  \centering
  \small
  \caption{Formal-mathematics benchmarks with the best publicly reported results
  as of August 2026. Competition-level proving is close to saturation;
  research-level tracks remain open, which motivates our setting. ``Aligned
  proofs'' marks benchmarks that ship NL--FL aligned proof material.}
  \label{tab:benchmarks}
  \begin{tabular}{@{}p{2.5cm}p{4.2cm}rcp{3.6cm}@{}}
    \toprule
    Benchmark & Task & Scale & Aligned proofs & SOTA (Aug.\ 2026) \\
    \midrule
    MiniF2F~\citep{minif2f2022} & Competition proving & 244 & \no &
      99.6\% (Seed-Prover) \\
    PutnamBench~\citep{putnambench2024} & Competition proving & 660 & \no &
      87.9\% (Seed-Prover 1.5) \\
    ProofNet~\citep{proofnet2023} & Undergrad.\ statement autoform. & 371 & \no &
      45.1\% ProofNet\# (GPT-4o) \\
    Gaokao-Formal~\citep{mathesis2025} & End-to-end NL proving & 488 & \no &
      18\% (Mathesis) \\
    RLMEval~\citep{rlmeval2025} & Research-level proving and proof autoform. &
      613 & \yes & 10.3\% \\
    FATE-H / FATE-X~\citep{fate2025} & Graduate+ algebra proving & 100 / 100 &
      \no & 3\% / 0\% (open provers) \\
    FormalQualBench~\citep{formalqualbench2026} & PhD-qual.\ theorem development
      & 23 & \no & 10/23 (MerLean-Prover) \\
    ProofFlow bench~\citep{proofflow2025} & Proof translation with step DAGs &
      184 & \yes & ProofScore 0.545 \\
    \bottomrule
  \end{tabular}
\end{table}


% \begin{table}[H]
%   \centering
%   \small
%   \caption{Formal-mathematics benchmarks (SOTA as reported by the sources).}
%   \label{tab:benchmarks}
%   \begin{tabular}{@{}p{2.1cm}p{4.5cm}rp{3.4cm}@{}}
%     \toprule
%     Benchmark & What it measures & Scale & SOTA \\
%     \midrule
%     MiniF2F~\citep{minif2f2022} & Formal proving from a Lean statement & 488 &
%       Goedel-Prover-V2-32B, 88.1\% pass@32 \\
%     RLMEval~\citep{rlmeval2025} & Research-level proving \& autoformalization &
%       613 & 10.3\% \\
%     ProofNet~\citep{proofnet2023} & Undergraduate statement autoformalization &
%       371 & 45.1\% (GPT-4o) \\
%     PutnamBench~\citep{putnambench2024} & Competition formal proving & 672 &
%       DeepSeek-Prover-V2, 49/658 \\
%     FormalMATH~\citep{formalmath2025} & Formal proving \& test-time scaling &
%       5{,}560 & DeepSeek-V2-671B, 28.31 \\
%     FormalQualBench & Graduate-level theorem development & 23 & OpenGauss, 8/23 \\
%     \bottomrule
%   \end{tabular}
% \end{table}

\section{Problem Formulation}
{\color{red}Given a natural-language pair $(NL_{\text{thm}}, NL_{\text{proof}})$, produce
$(FL_{\text{thm}}, FL_{\text{proof}})$ such that (i) the Lean kernel accepts
$FL_{\text{proof}}$ as a proof of $FL_{\text{thm}}$ using only permitted axioms
and without \texttt{sorry}; (ii) $FL_{\text{thm}}$ matches the intended
statement under the Comparator; and (iii) $FL_{\text{proof}}$ follows the
informal argument. To make (iii) measurable, we represent the informal proof as
a directed acyclic graph $D=(V,E)$ whose nodes are self-contained claims and
whose edges are dependencies, and we require the formal proof to take the form
of a \emph{skeleton} $S(D)$: one named lemma per node, closed individually, and
assembled into the main theorem. Proof faithfulness is the fraction of nodes in
$V$ with an aligned, kernel-closed counterpart in the final proof.

\paragraph{H1: Identifier grounding by construction.} After supervised
finetuning, the dominant failure of small formalizers is reference to
declarations that are absent from Mathlib (34.8\% of failures in our analysis in
Section~\ref{sec:results}; a frontier model shows about 50\% on ProofNet\# under
paraphrase~\citep{surfacesens2026}). We hypothesize that restricting decoding at
identifier positions to a trie over the declaration names of the pinned Mathlib
build, with retrieval supplying a ranked shortlist, removes this class at the
statement level by construction and raises Compile@1 and Faithful@1 beyond
retrieval-only grounding at equal sampling budgets. Refutation condition: gains
fail to exceed retrieval-only grounding, or forced near-miss names degrade
faithfulness. Prior work names this direction and leaves it
open~\citep{surfacesens2026}.

\paragraph{H2: Where structure-guided translation fails at research level.}
Structure is known to help: informal sketches seed formal proofs
in~\citet{dsp2023}, subgoal decomposition drives
DeepSeek-Prover-V2~\citep{deepseekproverv2}, lemma graphs outperform monolithic
translation in ProofFlow~\citep{proofflow2025}, and blueprint refinement carries
Goedel-Architect~\citep{goedelarchitect2026}. The open question is where the
residual failures live once grounding is controlled. We hypothesize that at
research level the majority of residual failures are node mis-statements, claims
that are false or unprovable as written, with leaf-closing failures second, and
that oracle structure contributes a larger marginal gain than oracle retrieval,
with a positive interaction between the two. Design: the six conditions of
Table~\ref{tab:factorial} on a pooled set of at least 300 problems from RLMEval,
FATE-H, and FormalQualBench, with every failure labeled into the five classes of
Table~\ref{tab:failures}. Refutation condition: leaf-closing failures dominate,
or retrieval matches structure within confidence intervals.

\paragraph{H3: Per-node semantic rewards for proof translation.} Dense syntactic
rewards help proving~\citep{processrl2026}, and semantic rewards help statement
formalization~\citep{formarl2025, mathesis2025, criticlean2025}. Their
combination for proof translation is untested. We hypothesize that, under
matched training and rollout budgets, GRPO with a per-node reward combining node
compilation, node provability by a frozen prover, and node-level semantic
consistency with the aligned informal step improves Verified@1 and proof
faithfulness over (a) an outcome-only compiler-and-Comparator reward and (b) a
per-node syntactic reward without the semantic term. Design: five arms
(whole-proof SFT, structured SFT, outcome-reward GRPO, per-node syntactic GRPO,
per-node semantic GRPO) with identical token budgets and identical policy
initialization. Refutation condition: the semantic term adds nothing over the
syntactic per-node arm, or the judge is gamed, inflating semantic scores without
gains in Verified@1.

Pass rates throughout are pass@1 estimated from $n{=}64$ samples per problem
with the unbiased estimator of~\citet{chen2021codex}, three seeds, and
stratified bootstrap 95\% confidence intervals; paired conditions are compared
with McNemar's test, and the per-problem outcome matrix is released.

\begin{table}[h]
    \centering
    \caption{Factorial design for H2. Grounding refers to the constrained
    identifier decoding of H1.}
    \label{tab:factorial}
    \small
    \begin{tabular}{l c c c c c c c}
        \toprule
        Condition & Formal stmt & Structure & Retrieval & Grounding &
          Compile@1 & Faithful@1 & Cost/prob. \\
        \midrule
        Direct generation      & \no  & \no  & \no  & \no  & & & \\
        Statement oracle       & \yes & \no  & \no  & \no  & & & \\
        Retrieval oracle       & \yes & \no  & \yes & \no  & & & \\
        Structure oracle       & \yes & \yes & \no  & \no  & & & \\
        Structure + retrieval  & \yes & \yes & \yes & \no  & & & \\
        Full (grounded)        & \yes & \yes & \yes & \yes & & & \\
        \bottomrule
    \end{tabular}
\end{table}

\begin{table}[h]
    \centering
    \caption{Failure decomposition. Every failed sample receives exactly one
    label; fractions are per condition.}
    \label{tab:failures}
    \small
    \begin{tabular}{l c c c c c}
        \toprule
        Condition & Parse & Unknown ident. & Node mis-stmt & Leaf unsolved &
          Assembly \\
        \midrule
        Direct generation      & & & & & \\
        Statement oracle       & & & & & \\
        Retrieval oracle       & & & & & \\
        Structure oracle       & & & & & \\
        Structure + retrieval  & & & & & \\
        Full (grounded)        & & & & & \\
        \bottomrule
    \end{tabular}
\end{table}


}
% How to translate natural-language mathematical problems and solutions into faithful verifiable statement-and-proof pairs, enabling scalable formal verification?

% Given a natural language theorem and proof pair $(NL_\text{theorem}, NL_\text{proof})$, we want to get the formal statement and proof $(FL_\text{theorem}, FL_\text{proof})$, which is verifiable by Lean and semantic faithful with the natural language pair. 

% \paragraph{Hypothesis 1: Proof-aware autoformalization enables more verified solutions. (confirm motivation)} Autoformalizing both the theorem and the reasoning in a natural-language solution will produce a higher rate of correctly verified solutions than autoformalizing the theorem alone and asking a prover to reconstruct the proof independently.
% \begin{table}[h]
%     \centering
%     \caption{Autoformalization enables more verified solutions}
%     \label{tab:h1}
%     \begin{tabular}{c|c c}
%     \toprule
%         Method & Pass@3 & Testing time \\
%         \midrule
%         statement formalization + formal reasoning & & \\
%         statement formalization + formal reasoning with NL context & & \\
%         joint theorem-and-solution autoformalization & & \\
%         \bottomrule
%     \end{tabular}
% \end{table}

% \paragraph{Hypothesis: Proof-structure recovery is the dominant remaining bottleneck (confirm bottleneck)} Given a correct Lean theorem statement, supplying oracle intermediate claims and dependency relations will improve solution's passrate more than supplying oracle library premises alone.

% \begin{table}[h]
%     \centering
%     \caption{Locate bottleneck}
%     \label{tab:h2}
%     \small
%     \begin{tabular}{c|c c c c c}
%         \toprule
%          Condition& formal statement & intermediate lemma statements & Mathlib declarations & Pass@3 & testing time \\
%         \midrule
%          Direct generation & \no & \no  & \no & & \\
%          Statement oracle&   \yes & \no & \no & & \\
%          Structure oracle & \yes & \yes& \no & & \\
%          Retrieval oracle & \yes & \no & \yes & & \\
%          Combined oracle & \yes & \yes & \yes & & \\
%          \bottomrule
%     \end{tabular}
% \end{table}

% \paragraph{Hypothesis 2: Step-level verifier feedback is more useful than final binary reward (inspire methods)}

Verifier-guided reinforcement learning applied to individual formalized claims
will improve end-to-end verification more than training from final compilation or proof-validity rewards alone. 

We compare whole-proof supervised
fine-tuning, structured supervised fine-tuning, final-reward reinforcement
learning, and step-level verifier-guided reinforcement learning under matched
training and rollout budgets.

\section{Method}
{\color{red}\label{sec:method}
Our pipeline has three stages, followed by training and a verification protocol.

\paragraph{Stage A: parsing and alignment.} A parsing model segments
$NL_{\text{proof}}$ into claims and dependency edges, restating quantifiers and
hypotheses so that every node is self-contained. The output is the graph
$D=(V,E)$. We follow the annotation conventions of ProofFlow's dependency
graphs~\citep{proofflow2025} and hand-validate a sample of parses per dataset.

\paragraph{Stage B: skeleton emission with grounded decoding.} The formalizer
emits a single Lean file: the target theorem signature is fixed verbatim, each
node of $D$ becomes a named lemma ending in \texttt{:= by sorry} with its
parents declared, and the file must elaborate under the pinned toolchain,
following the blueprint conventions of~\citet{goedelarchitect2026,
leanarchitect2026}. Grounding works as follows. We extract the full set of
declaration names from the pinned Mathlib build once, build a token-level trie
over them, and whenever the compiler reports an unknown identifier we re-decode
the offending span with logits masked to trie-valid continuations, placing a
retrieval-ranked shortlist of candidate declarations in context. The loop
repeats until the file elaborates or the repair budget is spent. We ablate a
fully constrained first pass against this repair-only variant; the repair-only
variant is simpler and already eliminates the failure class.

\paragraph{Stage C: node closing and assembly.} A frozen prover
(Goedel-Prover-V2-8B~\citep{goedelproverv2_2025}) attempts each node given only
its declared parents. The assembled proof is audited: Comparator statement
match, axiom whitelist checked via \texttt{\#print axioms}, a scan for
\texttt{sorry}, and rejection of \texttt{native\_decide} and top-level
\texttt{axiom} declarations, following the specification-based standard
of~\citet{formalqualbench2026} and the audit findings
of~\citet{faultsbench2026}.

\paragraph{Training.} The policy is the Stage A+B formalizer (7B); the prover
stays frozen so that improvements attribute to translation. We first run SFT on
Herald's 44k proof pairs~\citep{herald2025} restructured into skeleton format,
then GRPO with the reward
\[
  r \;=\; w_0\,\mathbf{1}[\text{skeleton elaborates}]
  \;+\; \frac{1}{|V|}\sum_{v\in V}\Big(a\,\mathbf{1}[v\text{ compiles}]
  \;+\; b\,\mathbf{1}[v\text{ closed}]
  \;+\; c\,\mathrm{sem}(v)\Big)
  \;+\; w_1\,\mathbf{1}[\text{final proof accepted}],
\]
where $\mathrm{sem}(v)$ scores the consistency of node $v$ with its aligned
informal step using a judge validated against 100 human-labeled node pairs
(agreement reported). Setting $c{=}0$ gives the syntactic arm of H3; collapsing
to the last term gives the outcome-only arm. Judge scores are capped and audited
on a held-out set to detect reward hacking.

\paragraph{Verification protocol.} All checking runs \texttt{lake env lean} on a
pinned Mathlib commit through LeanInteract or the Kimina Lean
Server~\citep{rlmeval2025, kimina_prover_2025}, with \texttt{maxHeartbeats} and
wall-clock caps reported. We report per-problem token and dollar costs,
following~\citet{goedelarchitect2026}.}
%We use the standard three-part pipeline and close the loop with the verifier:
% (1) an \textbf{autoformalizer} turns the problem into a Lean statement, (2) a
% \textbf{proof generator} writes a Lean proof, and (3) an \textbf{evaluator}
% checks it compiles and is faithful to the target theorem (no unsound tricks).
% The flow is:
% \[
%   \text{\LaTeX{} problem} \rightarrow \text{formulator} \rightarrow
%   \text{Lean proof} \xrightarrow{\texttt{lake env lean}} \text{validation log}
%   \rightarrow \text{verified proof},
% \]
% where failed candidates go back to the formulator with the compiler error log
% (and user feedback if needed).

% Following FormaRL~\citep{formarl2025}, the backbone is an RL loop with a compiler
% reward: $1$ if \texttt{Solution.lean} compiles under \texttt{lake env lean}
% \emph{and} passes the Lean comparator (the proof matches the theorem using only
% the permitted axioms), else $0$. To fight sparse reward, we add FormaRL's
% semantic-consistency reward and supervision mined from natural-language proofs.

\section{Experiment}
\paragraph{Metrics}
\begin{enumerate}
    \item \textbf{Compile@1}: the output is valid Lean that elaborates without
    errors on the pinned toolchain.
    \item \textbf{Verified@1}: the kernel accepts the proof, the Comparator
    confirms it proves the target statement, the axiom audit passes, and no
    \texttt{sorry} appears.
    \item \textbf{Statement faithfulness}: agreement of $FL_{\text{thm}}$ with
    $NL_{\text{thm}}$, scored by a judge validated on 100 human-labeled pairs
    with Cohen's $\kappa$ reported, following the validation practice
    of~\citet{beyondcompilation2026, mathesis2025}.
    \item \textbf{Proof faithfulness}: the fraction of nodes in $D$ with an
    aligned, kernel-closed counterpart in $FL_{\text{proof}}$.
\end{enumerate}
A solution can satisfy Compile@1 and still fail Verified@1, for example through
a \texttt{sorry}, a disallowed axiom, or a proof of a different statement, so we
report the metrics separately. All pass rates use the unbiased pass@$k$
estimator with $n{=}64$ samples~\citep{chen2021codex}, three seeds, and
stratified bootstrap 95\% confidence intervals.

% \paragraph{Metrics} \begin{enumerate}
%     \item Compilation success: a rollout counts as a success if it is valid Lean code that compiles without errors.
%     \item Proof validity: a solution is mathematically correct, meaning it proves the same statement as problem, uses only the allowed axioms, and is accepted by the Lean kernel. 
%     \item Semantic faithfulness: 
% \end{enumerate}

% A solution can compile successfully but still fail proof validity, for example if it contains a \texttt{sorry}, uses a disallowed axiom, or proves a different statement than intended. We report both metrics to separate basic code correctness from actual mathematical correctness.

\paragraph{Baselines} Statement-level: FormaRL~\citep{formarl2025} and
Kimina-Autoformalizer-7B~\citep{kimina_prover_2025}, each followed by the same
frozen prover. Proof-level: ProofBridge~\citep{proofbridge2025} and
ProofFlow~\citep{proofflow2025} with their released code at matched sampling
budgets. Prover-only reference: Goedel-Prover-V2-8B from the formal statement
alone~\citep{goedelproverv2_2025}. Agent reference: published FormalQualBench
results for OpenGauss and MerLean-Prover, reported for context and marked as
obtained under different budgets~\citep{opengauss2026, merlean2026}.
%\paragraph{Baselines}
%FormalRL, openGaussian, gordel-prover

\subsection{Results}\label{sec:results}
\paragraph{Failure Case}
We finetuned Llama-3-8B on the Kimina-Prover and Goedel-Prover proof
corpora~\citep{kimina_prover_2025, lin2025goedelprover} to obtain Llama-8B-SFT,
and compared it against DeepSeek-Prover-V2-7B under its chain-of-thought prompt
template~\citep{deepseekproverv2}, using the Comparator on
FormalQualBench~\citep{formalqualbench2026} with a pinned Mathlib commit.
% TODO: report n samples per problem, seeds, decoding parameters, Mathlib commit.
Every failed sample receives exactly one label from the taxonomy below, which
also defines the classes of Table~\ref{tab:failures}:
%We finetuned Llama-3-8B on~\cite{kimina_prover_2025, lin2025goedelprover} to obtain Llama-8B-SFT, and compared it against the state-of-the-art DeepSeek-%Prover-V2-7B~\cite{deepseekproverv2} using the standard Lean-prover evaluator Comparator on FormalQualBench. We analyze the typical failure modes of both models below:
\begin{enumerate}
    \item \textbf{Syntax/parse errors.} The base model fails predominantly (60.2\%) due to outputs that are not valid Lean at all: markdown headers, natural-language explanations, or informal theorem statements (e.g., a stray \texttt{\#\#\#} header or a \texttt{Theorem} line written in prose) leak directly into the \texttt{.lean} file, causing the parser to reject the proof before any semantic checking can occur. After finetuning, this failure mode drops sharply to 11.1\%, indicating that SFT is highly effective at teaching the model to respect Lean's surface syntax and output well-formed code.
    \item \textbf{Unknown identifiers.} The finetuned model's dominant failure
    mode shifts to referencing lemmas or definitions that are absent from
    Mathlib (34.8\%, up from 5.2\% in the base model). The pattern is consistent
    with the model having learned the style of Mathlib proofs, including
    plausible naming conventions, without reliable knowledge of the actual
    library contents. The same class appears at the frontier: a paraphrase probe
    reports unknown identifiers on about half of ProofNet\# items for a strong
    closed model~\citep{surfacesens2026}. This observation motivates H1.
    \item \textbf{Unsolved goals.} Both models leave a comparable fraction of
    goals unsolved (11.4\% base, 13.2\% SFT), which suggests that search
    failures respond weakly to finetuning and persist for both models. H2
    measures how much of this residue reflects mis-stated intermediate claims as
    opposed to genuine leaf difficulty.
  %  \item \textbf{Unknown identifiers.} In contrast, the finetuned model's dominant failure mode shifts to referencing lemmas or definitions that do not exist in Mathlib (34.8\%, up from 5.2\% in the base model). We attribute this to the SFT model having learned the \emph{style} of Mathlib-based proofs — plausible-sounding naming conventions, common lemma patterns — without having reliably memorized the actual library contents, leading it to hallucinate identifiers that look correct but are not defined.
%    \item \textbf{Unsolved goals.} Both models leave a comparable fraction of goals unsolved (11.4\% base vs.\ 13.2\% SFT), suggesting that genuine reasoning/search failures — as opposed to surface-level formatting issues — are only marginally affected by finetuning and remain a persistent bottleneck for both models.
\end{enumerate}

\paragraph{Ongoing Results}
FormalQualBench targets agent harnesses built on frontier models, where the
strongest published systems solve 10 of 23 problems~\citep{merlean2026};
end-to-end scores for the 7B class are therefore near zero and carry little
signal. We instead evaluate the effect of post-training stage by stage, on
statement formalization and on proving, across MiniF2F, ProofNet, and
FormalQualBench, whose difficulty increases in that order.
Table~\ref{tab:main} reports Compile@1 and Verified@1 for each cell.

\paragraph{Protocol note on finetuning RL-trained provers.}
DeepSeek-Prover-V2-7B ships as a reinforcement-learned checkpoint with distinct
chain-of-thought and non-chain-of-thought prompt modes~\citep{deepseekproverv2}.
In a preliminary run, supervised finetuning on top of this checkpoint reduced
MiniF2F accuracy from 70\% to 40\% at pass@8. Two effects explain drops of this
kind: a template mismatch that suppresses the model's chain-of-thought mode, and
the tendency of supervised finetuning to overwrite behavior acquired through
reinforcement learning~\citep{chu2025sft}. Our protocol therefore first
reproduces the published base numbers under the model's own template, applies
finetuning to base or instruct checkpoints, and reports the exact template,
decoding parameters, and Mathlib commit for every cell.

\begin{table}[h!]
    \centering
    \small
    \caption{Post-training influence on formalization and proving. Each cell
    reports Compile@1 (Verified@1) with 95\% bootstrap intervals. SFT rows start
    from base or instruct checkpoints; see the protocol note.}
    \begin{tabular}{l c c c}
        \toprule
        Model & MiniF2F & ProofNet & FormalQualBench \\
        \midrule
         Llama-3-8B & & & \\
         Llama-3-8B-SFT & & & \\
         Llama-3-8B-GRPO & & & \\
         DeepSeek-Prover-V2-7B (CoT) & & & \\
         DeepSeek-Prover-V2-7B-SFT & & & \\
         DeepSeek-Prover-V2-7B-GRPO & & & \\
         Kimina-Autoformalizer-7B & & & \\
         Goedel-Prover-V2-8B & & & \\
        \bottomrule
    \end{tabular}
    \label{tab:main}
\end{table}
\begin{figure}[h!]
    \centering
    \fbox{\rule{0pt}{4cm}\rule{0.45\linewidth}{0pt}}\hfill
    \fbox{\rule{0pt}{4cm}\rule{0.45\linewidth}{0pt}}
    \caption{Planned figure. Left: Compile@1 and Faithful@1 for the six
    conditions of Table~\ref{tab:factorial} with 95\% bootstrap intervals.
    Right: failure-class composition per condition as stacked bars, using the
    taxonomy of Table~\ref{tab:failures}.}
    \label{fig:factorial}
\end{figure}

% Since FormalQualBench is too complex for current models to complete the end-to-end math-proving task, we instead separately evaluate the influence of post-training on the base model's formalization and proving abilities. Specifically, we evaluate problem formalization ability on MiniF2F, ProofNet, and FormalQualBench, whose difficulty increases in that order.

% Table~\ref{tab:main} reports these results.

% \begin{table}[h!]
%     \centering
%     \caption{Post-training influence on formalization ability}
%     \begin{tabular}{c c c c}
%         \toprule
%         Model & MiniF2F & ProofNet & FormalQualBench\\
%         \midrule
%          Llama-3-8B & & & \\
%          Llama-3-8B-SFT & & & \\
%          Llama-3-8B-GRPO & & & \\
%          DeepSeek-Prover-V2-7B & & & \\
%          DeepSeek-Prover-V2-7B-SFT & & & \\
%          DeepSeek-Prover-V2-7B-GRPO & & & \\
%         \bottomrule
%     \end{tabular}
%     \label{tab:main}
% \end{table}

% \begin{figure}[h!]
%     \centering
%     \fbox{\rule{0pt}{4cm}\rule{0.5\linewidth}{0pt}}
%     \caption{Larger models achieve better performance.}
%     \label{fig:scale_size}
% \end{figure}
\section{Limitations}
Judge-based faithfulness scores can be gamed, so we validate the judge on human
labels and audit reward trajectories on a held-out set. The factorial covers
three factors and holds the prover and prompt format fixed, so interactions with
prover choice remain unmeasured. Research-level benchmarks are small, which is
why we pool sets and report intervals in place of single points. Results for 7B
models may order differently at larger scales.

\section{Reproducibility}
We release code, the per-problem outcome matrix, prompts, the Mathlib commit,
decoding parameters, seeds, and per-problem token and dollar costs. Verification
runs \texttt{lake env lean} with the Comparator, an axiom whitelist checked by
\texttt{\#print axioms}, a scan for \texttt{sorry}, and rejection of
\texttt{native\_decide} and top-level \texttt{axiom} declarations.


\bibliographystyle{abbrvnat}
\bibliography{ref}

\end{document}
